\input{template.tex}

\ifPDFTeX
  \renewcommand{\rmdefault}{cmr}
\fi

\def\SheetNo{3}
\setlength{\emergencystretch}{2em}

\begin{document}
\sheettitle{Множества и отображения.}

\begin{problems}

\item Верно ли, что множество всех функций $f\colon[0,1]\to\mathbb R$ имеет мощность континуума?

\item Бесконечное количество студентов написало контрольную, и каждый получил за неё какой-то балл --- число от $0$ до $5$, но никто, кроме практика, не знает баллы за контрольную. Практик им сообщил, что лишь конечное число студентов набрали балл, необходимый для зачёта, и предложил им сыграть в игру: студенты выстраиваются в шеренгу, и каждому на спине напишут его балл за контрольную. Каждый стоящий сзади видит баллы людей, стоящих перед ним, но не видит свой балл и баллы стоящих за ним. Затем с конца шеренги (с человека, который видит всех) практик начинает спрашивать, сколько баллов набрал студент (студент отвечает на ухо практику). Если студент угадает, то получит зачёт, иначе нет, даже если баллов на него хватает. Могут ли студенты сговориться так, чтобы зачёт по итогам игры получило бесконечное количество человек?

\item Верно ли, что если для некоторой функции $f$ и множеств $Y_1,Y_2$ выполняется $Y_1\cap Y_2=\varnothing$, то выполняется и $f^{-1}(Y_1)\cap f^{-1}(Y_2)=\varnothing$?

\item Придумайте какое-либо взаимно однозначное соответствие между разбиениями натурального числа на различные и на нечётные слагаемые.

\item Дана бесконечная последовательность цифр и нечётное число $d$, не делящееся на $5$. Докажите, что можно выбрать несколько цифр подряд, образующих число, делящееся на $d$.

\item После того как бесы, коих бесконечно много, проиграли Балде, им стало так стыдно, что они решили заняться физкультурой и записались в спортивные секции. Известно, что в каждую секцию записалось конечное число бесов и среди любого бесконечного набора бесов найдутся двое, занимающиеся в одной и той же секции. Докажите, что бесов-лентяев, которые ходят в конечное число секций, лишь конечное число.

\item Докажите, что следующие множества счётны:
\begin{parts}
  \item множество всех трёхэлементных подмножеств $\mathbb N$;
  \item множество всех конечных подмножеств $\mathbb N$.
\end{parts}

\item
\begin{parts}
  \item Существует ли континуальная цепочка попарно вложенных друг в друга подмножеств счётного множества?
  \item Может ли иметь континуальную мощность система подмножеств счётного множества, любые два из которых имеют конечное пересечение?
\end{parts}

\item Функция $f$ из $\mathbb N$ в $\mathbb N$ сопоставляет числу $n$ наибольший простой делитель числа $n$. Найдите полный прообраз множества чётных чисел.

\item Верно ли, что множество всех непрерывных функций из $\mathbb R$ в $\mathbb R$ континуально?

\item Покажите геометрически, что любые два отрезка равномощны и любой интервал равномощен прямой.

\item В клетки таблицы $3\times6$ записаны числа $1,2,\ldots,9$, каждое число по два раза. Будем называть таблицу хорошей, если в каждой строчке все числа различны и в первых трёх столбцах встречаются все числа от $1$ до $9$. Пусть $N$ --- количество хороших таблиц. Найдите степень вхождения $2$ в $N$.

\item Циклом $(a_1a_2\ldots a_n)$ называется перестановка (биекция на конечном множестве), при которой элемент $a_1$ переходит в $a_2$, в свою очередь, $a_2$ переходит в $a_3$ и так далее до $a_n$, который переходит в $a_1$. Остальные элементы множества переходят сами в себя (остаются на месте). При $n=2$ получается обмен двух элементов, или транспозиция; при $n=1$ получается тождественная перестановка. Покажите, что всякую перестановку можно представить как композицию нескольких непересекающихся циклов (более точно: не пересекаются множества тех элементов, которые циклы не оставляют на месте).

\item По кругу лежат карточки с числами от $1$ до $13$ именно в таком порядке. Изначально все карты лежат числами вверх. Вася может перевернуть карту, лежащую числом вверх, если карта, находящаяся на расстоянии $3$ от данной карты, лежит числом вверх. Например, $12$ можно перевернуть числом вниз, только если $9$ или $2$ лежат числом вверх. Сколько существует способов перевернуть $12$ карт числом вниз? В качестве ответа введите натуральное число.

\item Даны $1985$ множеств, каждое из которых состоит из $45$ элементов, причём объединение любых двух множеств содержит ровно $89$ элементов. Сколько элементов содержит объединение всех этих $1985$ множеств?

\item Докажите, что множество алгебраических чисел, то есть чисел, являющихся корнями многочленов с рациональными коэффициентами, счётно.

\item У конечного множества $A$ выбрано несколько подмножеств, каждое из которых содержит хотя бы $k^2$ элементов ($k\ge2$). Любые два множества пересекаются, причём не больше чем по $k-1$ элементу. Докажите, что элементы множества $A$ можно покрасить в два цвета так, чтобы каждое из выбранных подмножеств содержало элементы обоих цветов.

\item Пусть $X^Y=\{f\mid f\colon Y\to X\}$. Докажите, что существуют биекции между
\[
(X^Y)^Z\ \text{и}\ X^{Y\times Z},
\qquad
(X\times Y)^Z\ \text{и}\ X^Z\times Y^Z.
\]

\item Докажите, что квадрат равномощен отрезку.

\item Докажите, что множество всевозможных бесконечных последовательностей нулей и единиц имеет мощность континуума (то есть равномощно отрезку $[0,1]$).

\item
\begin{parts}
  \item На плоскости изображено бесконечное множество попарно не пересекающихся кругов. Докажите, что это множество счётно.
  \item Назовём восьмёркой фигуру на плоскости, состоящую из двух касающихся внешним образом окружностей. Докажите, что на плоскости можно расположить не более чем счётное множество попарно непересекающихся восьмёрок. Одна восьмёрка может располагаться внутри другой, не должны пересекаться только сами окружности.
\end{parts}

\item Может ли для некоторой функции $f$ и множеств $X_1,X_2$, таких что $X_1\cap X_2=\varnothing$, быть так, что $f(X_1)=f(X_2)$?

\item Строка длины $2023$ из $0$ и $1$ называется красной, если в ней есть $7$ единиц подряд. Строка длины $2024$ из $0$ и $1$ называется синей, если в ней есть $8$ одинаковых символов подряд. Каких строк больше, красных или синих, и во сколько раз?

\item Множество всех целых положительных чисел разбили на два непересекающихся подмножества --- $A$ и $B$. Докажите, что для любого целого $n>0$ найдутся такие целые $x>y>n$, что
\[
\{x,y,x+y\}\subseteq A
\]
или
\[
\{x,y,x+y\}\subseteq B.
\]

\item  $n \geq 2$ втшников решали $2^{n-1}$ задач. Оказалось, что для каждых двух задач есть ученик, который решил обе эти задачи, и ученик, который решил ровно одну из них. Докажите, что есть задача, которую решили все.


\end{problems}
\end{document}